% 电路与电子学III 课程笔记 —— 第3章（教材：孙立山《电路理论基础》第四版 第3章）
\chapter{电路定理}

\noindent
本章介绍电路理论中几个常用定理。\textbf{置换定理}适用于线性与非线性电路；\textbf{齐性定理}和\textbf{叠加定理}体现线性电路的本质，是线性方程的齐次性与可加性在电路中的体现；\textbf{戴维宁定理}和\textbf{诺顿定理}（合称等效电源定理）是化简线性含源一端口电路的有效方法；\textbf{特勒根定理}与基尔霍夫定律同样普遍适用，并可用以证明\textbf{互易定理}；对偶原理见 §2.8。

\section{置换定理}
\textbf{置换定理}：在任意线性和非线性电路中，若已知某一端口的电压和电流分别为 $U$ 和 $I$，则可用 $U_S=U$ 的电压源或 $I_S=I$ 的电流源来置换此一端口，而\textbf{不影响电路中其他部分}的电压和电流。

证明思路（图示法）：设一端口 $N_1$ 与 $N_2$ 的端口特性分别为 $U=f_1(I)$ 与 $I=f_2(U)$，把它们画在 $u-i$ 平面上，两曲线交点的坐标 $(U,I)$ 即端口的解。置换时用源电压等于该解的电压源（其图像是平行于 $i$ 轴的直线）代替 $N_2$，由于置换前后 $N_1$ 的结构与参数均未改变，$f_1(\cdot)$ 曲线不变，而端口电压仍为 $U$，故端口电流仍为 $I$——$N_1$ 内部各电压、电流均不受影响。用电流源置换的论证类似。

应用置换定理时必须注意两点：
\begin{enumerate}
  \item 置换后的电路必须\textbf{有唯一解}（实际电路与大多数电路模型都满足）；
  \item 除被置换部分发生变化外，\textbf{其余部分在置换前后必须保持完全相同}。
\end{enumerate}

由置换定理可得两个常用\textbf{推论}：
\begin{itemize}
  \item 若电路中某两点间电压为零，则可将量值为零的电压源接于该两点间，相当于把这两点\textbf{短路}，短路后不影响其他部分的解答；
  \item 若电路中某支路电流为零，则可将量值为零的电流源串接于该支路，相当于把该支路\textbf{断开}，断开后不影响其他部分的解答。
\end{itemize}
典型应用：电桥平衡时，桥臂上无电流的电阻可用短路线或断路线置换，从而把电桥化成串并联电路求等效电阻。

\corepoint{置换定理是“用一个数值已知的电源顶替一个已知端口”，它\textbf{不要求电路线性}（这是它与戴维宁定理最大的区别）；但置换的是“外部一端口”，被置换的那部分内部如何变化都无所谓，其余部分必须原封不动。}

\section{齐性定理和叠加定理}
齐性定理和叠加定理都是体现线性电路本质的重要定理。

\subsection{齐性定理}
\textbf{齐性定理}：在只有一个激励 $X$ 作用的线性电路中，设任一响应为 $Y$，记作 $Y=f(X)$，若将该激励乘以常数 $K$，则对应的响应 $Y'$ 也等于原来响应乘以同一常数，即
\keyformula{Y'=f(KX)=Kf(X)=KY}
证明思路：对含单一电压源的电路列回路电流方程，把源电压换成 $KU_S$ 后，方程各项乘以 $K$ 仍成立，两方程组系数相同，故解也成 $K$ 倍；电路中其他响应都是回路电流的线性组合，故同样变为 $K$ 倍。从数学上看，这是因为电路方程是\textbf{线性方程}。

对于线性直流电路，方程为线性代数方程，齐性定理可直观表述为：\textbf{若电路中只有一个激励，则响应与激励成正比}。即 $U_{\text{响应}}=\alpha_1U_S$ 或 $I_{\text{响应}}=\alpha_2I_S$，其中比例系数只与电路结构和参数有关。

\subsection{叠加定理}
\textbf{叠加定理}：在线性唯一解电路中，由几个独立电源共同作用产生的响应，等于各个独立电源\textbf{单独作用}时产生相应响应的\textbf{代数叠加}。

“单独作用”的含义：其余独立电源\textbf{置零}——电压源用短路代替、电流源用开路代替；受控源必须保留。

证明思路：把各独立电源共同作用时的电路方程，与每个独立电源单独作用时的电路方程逐项相加，所得方程与原方程完全相同，由解的唯一性即得叠加关系。能够这样做的根本原因是这些方程都是\textbf{系数相同}的线性方程，而各系数与某独立电源是否作用无关，只取决于电路结构、电阻及受控源参数。

注意：\textbf{功率不能直接叠加}。功率与源电压、源电流不是线性关系（$p=Ri^2$ 是关于 $i$ 的二次式），故
\begin{equation}
p=Ri^{2}=R\left(i'+i''\right)^{2}\neq R(i')^{2}+R(i'')^{2}
\end{equation}
正确做法是先用叠加定理求出电压、电流，再计算功率。

含受控源电路的叠加：受控源的参数影响方程的系数，故在\textbf{每个}独立电源单独作用的电路中都必须\textbf{保留受控源}，其控制量取该电路中的相应电压或电流。

对线性直流电路，由齐性定理与叠加定理可得统一结论：任一响应都是各独立电源的\textbf{线性组合}
\keyformula{Y=K_1X_1+K_2X_2+\cdots+K_mX_m}
其中 $X_j\ (j=1,2,\cdots,m)$ 为第 $j$ 个独立电源，$Y$ 为任一响应，$K_j$ 为常系数，取决于电路结构、电阻和受控源参数，\textbf{与独立电源的量值无关}。该式同时包含了齐性定理与叠加定理的含义。

\corepoint{用叠加定理最容易犯的三个错：（1）把受控源也“置零”（受控源必须保留）；（2）直接叠加功率；（3）各分响应叠加时忘记代数量的符号（方向）。另外“线性电路”是前提，含非线性元件时叠加定理不成立。}

\section{等效电源定理}
等效电源定理包括戴维宁定理和诺顿定理，是对线性含独立源一端口网络进行化简的重要定理。

\subsection{戴维宁定理}
\textbf{戴维宁定理}：线性含源一端口网络的对外作用，可以用一个\textbf{电压源串联电阻}的电路来等效代替；其中电压源的源电压等于此一端口网络的\textbf{开路电压}，而电阻等于此一端口网络内部各独立电源置零后所得无独立源一端口网络的\textbf{等效电阻}。所得电路称为\textbf{戴维宁等效电路}。

证明思路：在端口处外施电流源 $I$ 作激励，以端口电压 $U$ 为响应，由齐性定理与叠加定理
\begin{equation}
U=A_0I+\sum_{k}A_kX_k
\end{equation}
式中 $A_0$、$A_k$ 由电路结构与参数决定，与各独立源量值无关；$\sum A_kX_k$ 是与 $I$ 无关的确定电压，记作 $B_0$，于是
\begin{equation}
U=A_0I+B_0
\end{equation}
其中 $B_0$ 可这样确定：令端口所接电流源不作用（$I=0$，即端口开路），此时 $U=B_0$，故 $B_0=U_{oc}$（开路电压）；再令一端口内部各独立源不作用，此时 $U=A_0I$，可见 $A_0$ 是内部独立源全部置零时从端口看进去的等效电阻的负值，记 $A_0=-R_i$。于是
\keyformula{U=U_{oc}-R_iI}
这正是戴维宁等效电路（电压源 $U_{oc}$ 串联 $R_i$）的端口电压电流关系。

\subsection{诺顿定理}
戴维宁等效电路又可等效成电流源并联电导的电路，故含源一端口网络也可用\textbf{电流源并联电导}来等效。由两种含源支路的等效变换关系
\keyformula{I_{sc}=\frac{U_{oc}}{R_i},\qquad G_i=\frac{1}{R_i}}
式中 $I_{sc}$ 就是端口\textbf{短路}时戴维宁等效电路的端口电流，即含源一端口网络的\textbf{短路电流}。

\textbf{诺顿定理}：线性含源一端口网络的对外作用可以用一个电流源并联电导的电路来等效代替，其中电流源的源电流等于此一端口网络的短路电流，而电导等于内部各独立源置零后所得无独立源一端口网络的等效电导。所得电路称为\textbf{诺顿等效电路}。

由上述关系还得到求 $R_i$ 的第二种方法
\keyformula{R_i=\frac{U_{oc}}{I_{sc}}}
一般情况下含源一端口网络既可用戴维宁电路、也可用诺顿电路等效，视需要而定。特殊情况：
\begin{itemize}
  \item 若 $R_i=0$，则只能等效成戴维宁电路，并成为一个理想电压源；
  \item 若 $G_i=0$（$R_i\to\infty$），则只能等效成诺顿电路，并成为一个理想电流源。
\end{itemize}

\subsection{等效电阻的求法}
求 $R_i$ 的常用方法有三种，必须按电路是否含受控源来选择：

\begin{center}
\begin{tabular}{@{}lll@{}}
\toprule
电路情况 & 方法 & 说明 \\
\midrule
仅含电阻与独立源 & 串并联、$Y$-$\Delta$ 化简 & 独立源置零后直接化简 \\
仅含电阻与独立源 & $R_i=U_{oc}/I_{sc}$ & 分别求开路电压与短路电流 \\
含受控源 & 外加电源法 & 内部独立源置零，端口外加 $U$ 求 $I$，$R_i=U/I$ \\
\bottomrule
\end{tabular}
\end{center}

含受控源时\textbf{不能}只用电阻等效化简的方法求 $R_i$，因为受控源在“置零”后并不消失，它仍会影响端口特性。此时用外加电压源求端口电流（或外加电流源求端口电压）最稳妥。

工程意义：实际电路中常有某个特殊元件（如负载）经常变化而其他部分固定不变。等效电源定理为分析这类电路提供了简便方法——先求出不变部分的戴维宁或诺顿等效电路，把原电路化为简单电路再分析。

\corepoint{求戴维宁等效电路的三步：（1）断开待求支路，求开路电压 $U_{oc}$（注意极性，$U_{oc}$ 的“$+$”端就是等效电源的正极）；（2）内部独立源置零（电压源短路、电流源开路），求 $R_i$；（3）画等效电路再接回负载。含受控源时第（2）步必须用外加电源法。}

\section{特勒根定理}
\textbf{特勒根定理}是对所有集中参数电路都成立的定理，只需根据基尔霍夫定律即可证明；基尔霍夫定律只表述电路结构对电流、电压的约束，不涉及元件，因此特勒根定理的论证与应用均\textbf{与电路元件无关}。

设 $N$ 和 $\hat{N}$ 两个电路均有 $b$ 条支路、$n$ 个节点，且支路与节点的连接关系相同（称为\textbf{结构相同}的电路）。在 $N$ 和 $\hat{N}$ 中使对应支路、节点编号分别相同，并设对应支路具有相同的关联参考方向。

\textbf{特勒根定理}：电路 $N$ 中各支路电压与电路 $\hat{N}$ 中对应支路电流的乘积之和等于零，即
\keyformula{\sum_{k=1}^{b}u_ki_k'=0,\qquad \sum_{k=1}^{b}u_k'i_k=0}
证明：设支路 $k$ 接在节点 $a$、$b$ 之间，用 $\varphi_a$、$\varphi_b$ 表示 $N$ 的节点电压，则 $u_k=\varphi_a-\varphi_b$，代入求和式并按节点重新归并，每个节点前的系数是该节点流出电流之和——由 KCL 为零，故总和为零。

特勒根定理的特点是在两个电路 $N$ 与 $\hat{N}$ 之间建立了电压与电流的量值关系，\textbf{没有物理意义}。若将定理用于同一电路（即 $\hat{N}$ 也是 $N$），得到
\keyformula{\sum_{k=1}^{b}u_ki_k=0}
它表明：在任一时刻，一个电路中各支路吸收功率的代数和等于零，这就是电路的\textbf{功率守恒定理}。其涵义是：电路工作时某些支路吸收功率，同时另一些支路发出功率（吸收功率为负值），二者必须平衡。

当 $N$ 与 $\hat{N}$ 元件参数不同时，$\sum u_ki_k'$ 虽有功率的量纲，但不代表任何功率，只能称为\textbf{似功率}，故特勒根定理又称\textbf{似功率守恒定理}。

\corepoint{特勒根定理只要求“结构相同”，对元件性质毫无要求——可以是线性/非线性、时变/非时变，这正是它比叠加定理、戴维宁定理更普遍的原因。用它证互易定理时，取 $\hat{N}$ 为“把激励搬到另一端口”的电路即可。}

\section{互易定理}
有些电路（例如仅含线性二端电阻和独立源的电路），其回路电流方程与节点电压方程的系数矩阵都关于主对角线\textbf{对称}。具有这种性质的电路属于\textbf{互易性电路}，服从互易定理。互易定理有三种表述：

\textbf{第一种形式}：对于由一个独立电压源和若干线性二端电阻组成的电路，当此电压源在某一端口 $A$ 作用时，在另一端口 $B$ 产生的\textbf{短路电流}，等于把此电压源移到端口 $B$ 作用而在端口 $A$ 所产生的短路电流。

\textbf{第二种形式}：对于由一个独立电流源和若干线性二端电阻组成的电路，当此电流源在某一端口 $A$ 作用时，在另一端口 $B$ 产生的\textbf{开路电压}，等于把此电流源移到端口 $B$ 作用而在端口 $A$ 所产生的开路电压。

\textbf{第三种形式}：对于由一个独立电源和若干线性二端电阻组成的电路，如果在数值上 $I_S$ 与 $U_S$ 相等，则 $U_2$ 与 $I_1$ 在数值上也相等（$I_S$ 与 $I_1$、$U_S$ 与 $U_2$ 分别取同样单位）。

证明（第一种形式）：把两种情形的电路分别视为 $N$ 和 $\hat{N}$，二者结构相同且方框内只含二端电阻，服从特勒根定理。设端口 $A$、$B$ 所接支路分别为支路 1 和 2，方框内支路编号为 3 到 $b$，则
\begin{equation}
u_1i_1'+u_2i_2'+\sum_{k=3}^{b}u_ki_k'=0,
\qquad
u_1'i_1+u_2'i_2+\sum_{k=3}^{b}u_k'i_k=0
\end{equation}
方框内二端电阻上 $u_k=R_ki_k$、$u_k'=R_ki_k'$，两式的求和项相等，相减得
\begin{equation}
u_1i_1'+u_2i_2'=u_1'i_1+u_2'i_2
\end{equation}
在第一种形式中，情形 (a) 里 $u_2=0$（端口 $B$ 短路）、$u_1=U_S$；情形 (b) 里 $u_1'=0$（端口 $A$ 短路）、$u_2'=U_S$，代入即得 $U_Si_1'=U_Si_2$，从而两短路电流相等。第二、第三种形式可用对偶原理或类似方法得到。

综观互易定理的三种形式，有一个共性：\textbf{当激励不作用时，互易前后的两个电路是完全一样的}。注意这一共性，便于正确应用互易定理。

\corepoint{互易定理只适用于\textbf{互易性电路}（仅含线性二端电阻与独立源）。含受控源的电路一般不是互易电路，系数矩阵不再对称，互易定理失效。另外，“短路电流互易”与“开路电压互易”不要混用：电压源激励对应短路电流响应，电流源激励对应开路电压响应。}

\section{对偶原理（补充：对偶电路的形成）}
对偶因素与对偶原理的基本内容见 §2.8。这里补充教材 3.6 节给出的两个要点。

\textbf{对偶电路的直接画法}（以平面电路为例）：在每一个网孔内标出一个节点，作为对偶电路的独立节点；在外网孔外标出一个节点，作为对偶电路的非独立节点（参考节点）。把所标出的各节点用虚线互相连接，每条连线只通过一个元件，把该元件换成\textbf{对偶元件}（电阻 $\leftrightarrow$ 电导、电压源 $\leftrightarrow$ 电流源、电感 $\leftrightarrow$ 电容），即得对偶电路相应支路的元件。这样得到的电路，其节点电压方程与原电路的网孔电流方程互为对偶。

\textbf{工程应用}：利用对偶电路可以简化电路设计，例如把低通滤波器的设计结果直接用于高通滤波器的设计（见第 7 章滤波器简介）。

\begin{center}
\begin{tabular}{@{}ll@{}}
\toprule
原表述 & 对偶表述 \\
\midrule
KCL：节点上 $\sum i=0$ & KVL：回路上 $\sum u=0$ \\
电阻串联总电阻等于各电阻之和 & 电导并联总电导等于各电导之和 \\
电阻串联分压规律 & 电导并联分流规律 \\
戴维宁定理 & 诺顿定理 \\
互易定理第一种形式 & 互易定理第二种形式 \\
自阻、互阻 & 自导、互导 \\
星形（$Y$）联结 & 三角形（$\Delta$）联结 \\
\bottomrule
\end{tabular}
\end{center}

\section*{本章定理一览}
\addcontentsline{toc}{section}{本章定理一览}
\begin{center}
\begin{tabular}{@{}p{0.2\textwidth}p{0.36\textwidth}p{0.32\textwidth}@{}}
\toprule
定理 & 内容 & 适用条件 \\
\midrule
置换定理 & 端口 $(U,I)$ 已知时可用 $U_S=U$ 或 $I_S=I$ 置换该端口 & 线性或非线性电路，置换后有唯一解 \\
齐性定理 & 单一激励乘 $K$，响应也乘 $K$ & 线性电路 \\
叠加定理 & 多源响应等于各源单独作用响应之代数和 & 线性唯一解电路，受控源须保留，功率不叠加 \\
戴维宁定理 & 等效为 $U_{oc}$ 串联 $R_i$ & 线性含源一端口 \\
诺顿定理 & 等效为 $I_{sc}$ 并联 $G_i$ & 线性含源一端口 \\
特勒根定理 & $\sum u_ki_k'=0$，$\sum u_k'i_k=0$ & 两个结构相同的集中参数电路 \\
功率守恒 & 同一电路 $\sum u_ki_k=0$ & 集中参数电路 \\
互易定理 & 激励与响应互换位置，响应之比等于激励之比 & 互易性电路（线性二端电阻＋独立源） \\
对偶原理 & 对偶因素置换后命题仍成立 & 对偶的电路与方程 \\
\bottomrule
\end{tabular}
\end{center}

\endinput
